\section{归约的运算结构与并行代价}
\label{l15:sec:structure}

\subsection{二元运算、单位元与结果的含义}

归约（reduction）把一组输入通过同一种二元运算合并成一个结果。设输入包含 $N$ 个值，依次为 $x_0,x_1,\ldots,x_{N-1}$，二元运算记为 $\circ$，则结果为 $r=x_0\circ x_1\circ\cdots\circ x_{N-1}$。求和、求积、最小值和最大值都具有这种结构。与归约相关的前缀扫描（prefix scan）会输出各个前缀的累计结果，例如求和扫描在位置 $k$ 输出 $\sum_{i=0}^{k}x_i$。两者共享规则化的合并结构，归约只保留最终的一个值。

并行化需要明确运算允许怎样重排。\textbf{结合律} $(a\circ b)\circ c=a\circ(b\circ c)$ 允许改变括号结构，从而把一条依赖链改写成一棵树；\textbf{交换律} $a\circ b=b\circ a$ 允许交换元素顺序。后面按两半配对、再按线程束合并的求和实现都会改变分组，某些映射还会改变合并次序，因此采用满足这两条性质的理想运算模型。一般而言，保持原有顺序的归约树只依赖结合律；要允许任意顺序重排，还需要交换律。

\textbf{单位元} $e$ 满足 $e\circ a=a\circ e=a$。它既用于初始化累计量，也用于把不足一个完整数据段的位置补齐。常用运算与单位元如表\ref{l15:tab:identity} 所示。对于浮点最小值和最大值，这里只讨论普通有序数值，特殊非数值的处理需要另行指定。

\begin{table}[htbp]
\centering
\caption{常见归约运算及其单位元}
\label{l15:tab:identity}
\begin{tabularx}{\linewidth}{@{}lYYY@{}}
\toprule
运算 & 二元合并 & 单位元 & 补齐位置的含义\\
\midrule
求和 & $a+b$ & $0$ & 不增加总和\\
求积 & $ab$ & $1$ & 不改变已有乘积\\
最小值 & $\min(a,b)$ & $+\infty$ 或类型上界 & 不成为更小的候选\\
最大值 & $\max(a,b)$ & $-\infty$ 或类型下界 & 不成为更大的候选\\
\bottomrule
\end{tabularx}
\end{table}

顺序归约可以先令 $a=e$，再对每个输入更新 $a\leftarrow a\circ x_i$。这需要 $N$ 次更新；若 $N\geq1$，也可以直接用 $x_0$ 初始化，只合并余下的 $N-1$ 个值。两种写法的时间复杂度均为 $\Theta(N)$，后者便于与树上必需的 $N-1$ 次二元合并比较。空输入则按单位元约定返回 $e$。

整数求和也需要明确数值范围。数学整数的结合律不意味着固定宽度类型可以忽略溢出：有符号整数的某个中间部分和溢出时，C++ 运算不再具有预期语义，即使最终数学结果可表示也不足以保证安全。使用无符号类型时，加法按位宽取模；只有所有累计值均在计数范围内，模运算结果才等于通常的整数总和。

对实数加法成立的结合律，在有限精度浮点运算中通常只能近似成立。以单精度数值为例，$(10^{20}+(-10^{20}))+1$ 可以得到 $1$，而 $10^{20}+((-10^{20})+1)$ 可以得到 $0$，因为小量 $1$ 在与大幅值负数相加时被舍入掉。归约树改变括号结构，块间原子累加还可能改变执行次序，因此正确的并行求和也可能与顺序求和产生数值差异；存在严重消去时，这种差异可能明显增大。验证时应依据数据尺度比较绝对误差和相对误差；若任务要求固定的浮点结果，则还需要固定跨块合并次序及运算实现。

\subsection{从顺序依赖链到归约树}

以下并行代价讨论取 $N\geq2$。对于 $N=2^L$ 个输入，平衡二叉归约树在第一级完成 $N/2$ 次两两合并，第二级完成 $N/4$ 次合并，直到最后一级只合并两个部分结果。同一级的各个操作互不依赖，后一层只等待自己的两个子结果。图\ref{l15:fig:tree} 用八个数的求和说明这一过程。

\begin{figure}[htbp]
\centering
\begin{tikzpicture}[x=1.55cm,y=1.04cm,
  every node/.style={font=\small},
  value/.style={draw=UIUCBlue,fill=SoftBlue,minimum width=7mm,minimum height=6mm},
  edge/.style={-{Stealth[length=1.6mm]},draw=UIUCBlue}]
\foreach \i/\v in {0/4,1/7,2/2,3/3,4/8,5/5,6/9,7/6}
  \node[value] (x\i) at (\i,0) {$\v$};
\foreach \i/\v/\pos in {0/11/0.5,1/5/2.5,2/13/4.5,3/15/6.5}
  \node[value] (a\i) at (\pos,-1) {$\v$};
\foreach \a/\b/\c in {0/1/0,2/3/1,4/5/2,6/7/3}{
  \draw[edge] (x\a)--(a\c);\draw[edge] (x\b)--(a\c);}
\node[value] (b0) at (1.5,-2) {$16$};
\node[value] (b1) at (5.5,-2) {$28$};
\draw[edge] (a0)--(b0);\draw[edge] (a1)--(b0);
\draw[edge] (a2)--(b1);\draw[edge] (a3)--(b1);
\node[value] (c) at (3.5,-3) {$44$};
\draw[edge] (b0)--(c);\draw[edge] (b1)--(c);
\end{tikzpicture}
\caption{八个输入通过三级合并得到总和。每层分别执行四次、两次和一次加法。}
\label{l15:fig:tree}
\end{figure}

用\textbf{工作量}（work）$W(N)$ 表示所有线程执行的必要二元合并总数，用\textbf{跨度}（span）$D(N)$ 表示最长依赖路径上的合并数。在单位合并时间的抽象模型中，归约树满足
\begin{equation}
W(N)=\sum_{k=1}^{L}\frac{N}{2^k}=N-1,
\qquad D(N)=L=\log_2N.
\label{l15:eq:work-span}
\end{equation}
工作量与优化过的顺序算法相同，因此这棵树是工作高效的（work-efficient）。这里比较的是必要合并数；索引计算、数据搬移、同步和后面线程束代码中的冗余通道运算仍然具有实际开销。

若 $N$ 不是二的幂，可以补单位元到 $\widehat N=2^{\lceil\log_2N\rceil}$。规则树的深度为 $\lceil\log_2N\rceil$，工作量为 $\widehat N-1<2N-1$；若只合并实际存在的值，仍只需要 $N-1$ 次合并。后面的分块实现通过边界判断补零，无需真的建立一份全局填充数组。

平均可用并行度为 $W/D$。当 $N$ 是二的幂时，它等于 $(N-1)/\log_2N$，峰值出现在第一级，为 $N/2$。例如，对 $N=1{,}000{,}000$ 个实际输入建立平衡归约树时，深度为 $\lceil\log_2N\rceil=20$，工作量为 $N-1$，平均并行度约为 $50{,}000$，第一级包含 $500{,}000$ 次独立合并。这些数量说明，要把树深度直接兑现成运行时间，需要非常多的执行资源。

若硬件同一时刻最多执行 $P$ 次合并，则 $T_P\geq\max\{W/P,D\}$；忽略通信等开销的理想调度时间量级为 $\Theta(W/P+D)$。因此，$\Theta(\log N)$ 描述依赖深度，并不能单独给出实际图形处理器（GPU）的运行时间。树越接近根，活跃操作越少，后期必然出现并行度收缩。一个核函数启动后线程块的线程数固定，停止做加法的线程也不会自动变成新线程去处理下一批数据。

\subsection{原子累加与局部合并}

最直接的并行求和方式，是让每个线程读取一个输入，再用原子操作（atomic operation）把它加到同一个全局输出上。原子加法把一次读、加和写作为不可分割的更新，保证同时到达的线程不会因覆盖彼此的写入而丢失贡献。

\noindent\begin{minipage}{\linewidth}
\begin{lstlisting}[caption={逐元素原子求和：同一个输出地址接收全部更新},label={l15:lst:code11}]
__global__ void reduce_atomic(const float* input,
                              size_t n, float* sum) {
    size_t i = size_t(blockIdx.x) * blockDim.x
             + threadIdx.x;
    if (i < n) {
        atomicAdd(sum, input[i]);
    }
}
\end{lstlisting}
\end{minipage}

这段核函数要求输出地址有效，并且在核函数执行前将 \mbox{\code{*sum}} 设为 $0$。输出不能与仍待读取的输入区域重叠。即使所有输入读取都能并行，同一地址上的 $N$ 次原子更新仍然相互竞争；这些更新必须以某种顺序生效，难以获得归约树同一级那样的独立合并能力。原子操作保证更新正确，并不会自动消除这个共享地址形成的瓶颈。

一种直接的改进是先让每个线程块完成局部归约，只由该块的一个线程原子累加块和。若每块处理 $2B$ 个元素，块数为 $G=\lceil N/(2B)\rceil$，则全局原子更新数从 $N$ 降为 $G$。例如 $N=1{,}048{,}576$、$B=256$ 时，只需 $2048$ 次块结果更新。局部树负责把大量独立加法分散到各块中，最终的原子操作只连接这些块的结果。

另一种块间合并方法是先把 $G$ 个块和写入独立数组，再启动后续归约核函数，直到只剩一个值。同一执行流中的后续核函数会在先前核函数完成后使用这些部分和。这两种方案共享相同的块内算法，区别集中在最终结果如何跨块汇合；后面将分别给出它们的使用条件。
